\section{Sibling-Guided Credit Distillation}
\label{sec:method}

\subsection{Overview}
\method (Figure~\ref{fig:method}) augments a group-relative policy-gradient
backbone with a distillation-derived credit signal. Dynamic sampling produces
mixed successful/failed sibling rollouts; an external LLM turns them into a
stepwise credit reference; the policy is scored under clean and credit-reference
contexts; and detached teacher/student divergence reshapes token advantages
inside the policy-gradient update. The inference-time model is the ordinary student policy: no external LLM,
credit reference, or sibling evidence appears at inference.

\subsection{Mixed sibling evidence}
\label{sec:siblings}
For a task $x$, the rollout worker samples a sibling group
$Y_x=\{y_{x,1},\ldots,y_{x,n}\}$ under the current policy, with verifier outcome
$R_{x,k}\in\{0,1\}$. A group is useful when it is \emph{mixed}:
\begin{equation}
  0 < \sum_{k=1}^{n} R_{x,k} < n .
\end{equation}
Mixed groups contain verified successful routes and failed alternatives under
the same task prompt, making dynamic sampling SGCD's sibling evidence source.
Following DAPO-style dynamic sampling \citep{dapo2025}, all-success and
all-failure candidate groups are discarded and replacement task groups are
sampled. The reported recipe refills toward eight complete mixed groups for at
most ten generation batches; if the target remains unmet, it skips that
optimizer update and continues training. Homogeneous groups never enter an SGCD
policy update.

\paragraph{Credit-reference construction.} For each mixed group, a strong
external LLM converts policy-visible sibling traces plus binary verifier
outcomes into a training-only stepwise credit reference $z_x$: reusable state
checks, useful actions, failed
branches, and likely deviation points. We rebuild this reference from fresh
on-policy siblings rather than a persistent skill bank (prompt details:
\suppref{sec:appendix}).

\paragraph{Leakage boundary.} The credit reference is never shown to the
inference-time student. During construction we retain only rollout-visible
traces and the binary verifier outcome; all other evaluator-private fields,
gold actions, sensitive values, and answer-bearing literals are removed. The
reference therefore abstracts reusable credit evidence rather than hidden
answers.

\subsection{Dense teacher/student divergence}
\label{sec:dense}
For a rollout token position $t$, let $p_{x,k,t}^{K}$ be the current policy
scored on the clean rollout prefix, and let $q_{x,k,t}^{K}$ use the same policy
weights on the same observed prefix with credit reference $z_x$ inserted as
teacher context. Both distributions are projected onto the same local support.
The reported runs use $K{=}100$. The support
contains the teacher top-$K$ tokens, the observed rollout token, and a tail
bucket that preserves remaining mass. We primarily use reverse-KL,
\begin{equation}
  d_{x,k,t}
  = D_{\mathrm{KL}}\!\left(
      p_{x,k,t}^{K} \;\middle\|\; \mathrm{sg}[q_{x,k,t}^{K}]
    \right),
  \label{eq:rkl}
\end{equation}
The reported SGCD results use reverse-KL. Equation~\ref{eq:rkl} stops gradients
through $q$; the complete divergence/entropy feature is detached from the actor
through $u$ below. Thus the divergence and teacher entropy are consumed as
credit features, not as a direct actor-gradient distillation loss.

\subsection{Divergence-driven credit weighting}
\label{sec:segments}
Long-horizon traces contain large stretches of boilerplate reasoning, setup, and
tool output. SGCD uses token-level divergence for \emph{bounded credit
emphasis}: it adds local emphasis to the broadcast trajectory advantage without
claiming causal ground-truth credit or conserving advantage mass. Within each
response, our deterministic divergence/entropy segmenter uses
$\tilde d_{x,k,t}=(d_{x,k,t}-d_{\min})/(d_{\max}-d_{\min}+\epsilon_{\mathrm{norm}})$
with response extrema $d_{\min},d_{\max}$ and $\epsilon_{\mathrm{norm}}=0.1$.
A segment starts when $\tilde d_{x,k,t}>\lambda_{\mathrm{KL}}=0.15$ and extends
until entropy exceeds $\lambda_H=1.5$ times its onset value, subject to the
$0.2T$ cap in \suppref{sec:appendix}. Triggered-span tokens inherit the
onset saliency, while non-triggered tokens use their local $\tilde d_{x,k,t}$.
The saliency $s_{x,k,t}\in[0,1]$ marks confident, localized disagreement.
API/action spans are a safety/reporting slice, not the only segmentation rule.
The optimization mask includes only assistant-generated response tokens;
prompts, tool observations, simulator returns, and the credit-reference prefix
are excluded from the objective below.
Let $\mathcal{M}$ denote the resulting set of assistant-response token
positions.

\subsection{Credit-weighted policy gradient}
\label{sec:credit}
Let $A_{x,k}$ be the group-relative trajectory advantage and $\rho_{x,k,t}(\theta)$
the usual policy ratio. SGCD replaces the broadcast advantage with a
token-weighted advantage
\begin{equation}
  \begin{aligned}
  \widehat A_{x,k,t} &= A_{x,k} \, W_{x,k,t},\\
  s_{x,k,t} &= g(u_{x,k,t}),\\
  u_{x,k,t} &= \mathrm{sg}\!\left[d_{x,k,t},
  H(q_{x,k,t}^{K})\right],
  \end{aligned}
  \label{eq:weighted_adv}
\end{equation}
where $g$ is the deterministic divergence-spike and entropy-termination
segmenter from \S\ref{sec:segments}. In our implementation, the credit
multiplier is
\begin{equation}
  W_{x,k,t}
  = \mathrm{clip}\!\left(1+\gamma\,s_{x,k,t},
  1, c\right),
  \label{eq:weight_layer}
\end{equation}
where $\mathrm{clip}(a,\ell,u)=\min(\max(a,\ell),u)$, $c>1$, and
$\gamma\ge0$; reported runs use $\gamma{=}1$ and $c{=}2$. Low-saliency tokens
have $W\approx1$; high-saliency tokens amplify the existing advantage magnitude
without flipping its sign. The stop-gradient detaches feature construction from
actor and teacher scoring, and the multiplier is constant with respect to actor
parameters: there is no learned credit-weight network and no gradient through
$d_{x,k,t}$, $H(q^K)$, or a teacher-matching KL/JSD objective. Failed rollouts
receive larger negative advantage on high-saliency tokens; successful rollouts
receive concentrated positive credit; low-saliency tokens recover ordinary GRPO. The
actor loss is the standard
clipped surrogate with $\widehat A_{x,k,t}$. Let
$\bar\rho_{x,k,t}=\mathrm{clip}(\rho_{x,k,t},1-\epsilon,1+\epsilon)$; we use
symmetric $\epsilon{=}0.2$ (bounds $[0.8,1.2]$) as in PPO \citep{ppo2017}:
\begin{equation}
  \begin{aligned}
  \mathcal{L}_{\text{SGCD}}
  ={}& - \sum_{(x,k,t)\in\mathcal{M}} \\
     &\min\!\left(\rho_{x,k,t}\widehat A_{x,k,t},
      \bar\rho_{x,k,t}\widehat A_{x,k,t}\right).
  \end{aligned}
  \label{eq:sgcd_loss}
\end{equation}
This equation shows SGCD's credit-weighted PG component. Additional
regularization is recipe-specific and outside this notation. The distillation
pathway contributes no
teacher-matching actor-gradient term. The dense
teacher/student comparison and credit-reference construction determine
\emph{where} the verified advantage receives extra emphasis, not \emph{what}
token distribution the actor must imitate. The Vanilla GRPO comparators retain
their benchmark reference-KL regularizers: coefficient $0.05$ on
$\tau^3$-airline and $0.04$ on AppWorld, the latter following DeepSeekMath
\citep{deepseekmath2024}. SGCD omits reference-KL and instead uses detached
dense divergence as a replacement credit signal. This signal is not an
equivalent KL penalty because it has no direct teacher-matching actor-gradient
path. Finally, $\lambda_{\mathrm{KL}}=0.15$ in \S\ref{sec:segments} is a
normalized-divergence segment-onset threshold, not a loss coefficient.

\subsection{Why this removes the direct-SD failure path}
\label{sec:safety}
Earlier direct auxiliary variants have the form
$\mathcal{L}_{\text{PG}}+\lambda\mathcal{L}_{\text{SD}}$: if the teacher target
is wrong, privileged, or unreproducible, the auxiliary creates a second actor
gradient that can oppose the verifier. SGCD removes that path. The credit
reference and dense RKL are detached features; they identify where the verified
advantage receives bounded emphasis, while SGCD adds no distillation-derived
actor-gradient path beyond the reward-grounded PG term. This operationalizes the diagnosis in
\S\ref{sec:diagnosis}: use distillation for credit weighting, not to overrule
policy gradient.
